% ==========================================================================
% GENERATED by scripts/split_paper.py. Do not edit this file.
%
% Source: imported_papers/Neural_Computation__CNA/NECO-D-24-00304-Source.tex
% Part 3: Cooperative Network Architecture
%
% Edits belong upstream, in the published paper. To change what the
% chapter includes, change imports/cna.toml and re-run
% the splitter.
% ==========================================================================
\section{Cooperative Network Architecture}\label{sec:coop_net_architecture}

\begin{figure}[ht!]
  \begin{center}
    \includegraphics[width=\textwidth]{NECO-D-24-00304-Figure.2}
  \end{center}
  \caption{\textbf{Dynamics within the proposed CNA}: The input image is fed into stage $S1$ to obtain feature activations. The feature activations are then processed by stage $S2$, together with activations from neurons within the same layer. %In part A, two distinct matrices are used to model the forward connections from $S1$ to $S2$ and the recurrent connections within $S2$.
  %In part B, the matrices and input signals are concatenated, resulting in one bigger weight matrix. }
  }\label{fig:system_dynamics}
\end{figure}

In this section, we present a concrete instantiation of the proposed framework, specifically applied to the visual domain -- although the framework is general and may be extended to other sensory modalities in the future. We refer to this implementation as the Cooperative Network Architecture (CNA).
To realize our framework, we adopt a standard convolutional neural network (CNN) architecture~\citep{lecun_backpropagation_1989}.

The CNA comprises two processing stages, depicted as yellow squares in Figure~\ref{fig:system_dynamics}. \emph{Stage 1 ($S1$)} performs fixed (i.e., not learned) feature extraction directly from the image input. This fixed feature stage is used here for simplicity and to facilitate clearer diagnostics of the subsequent results by combining well-characterized, predefined features. However, these features could also be learned in future work, for example, through biologically inspired Hebbian learning schemes that have been shown to induce convolutional, Gabor-like filters \citep[e.g.,][]{miconi_hebbian_2021}. \emph{Stage 2 ($S2$)} implements the core principles of the described computational framework by learning net fragments.

The network dynamics unfold over discrete time steps $t = [0, \dots, T]$.
Since the feature extraction stage ($S1$) employs four fixed filters, it yields a time-invariant activation pattern for a given (fixed) input.
In contrast, the second stage ($S2$) exhibits dynamic temporal evolution, where internal representations are iteratively refined at each time step ($t \to t+1$).
Specifically, we use a progressively increasing global attenuation signal that only allows neurons with a high pre-synaptic activity to remain active.
The following sections discuss these two processing stages.

\subsection{Stage 1: Feature Extraction}
Stage 1 is designed to model the sensory input by extracting initial neural activations from the raw image data. This extraction is essential due to the continuous nature of typical image inputs (e.g., grayscale or RGB), whereas the computational framework employed in this work operates on binary activations. Consequently, this stage serves as a preprocessing step, enabling the transformation of continuous-valued images into a suitable binary representation.

The input image is denoted as $\boldsymbol{x} \in \mathbb{R}^{C, H, W}$ (illustrated by the gray box in Figure \ref{fig:system_dynamics}), where $C$ indicates the number of image channels ($1$ for grayscale, $3$ for RGB), while $H,  W$ stand for the height and width of the image.

\paragraph{Activation Computation.}
To extract initial activation patterns, we use four fixed, non-trainable filters (details see Appendix \ref{sec:features}). These filters are implemented as standard convolutional kernels, used to obtain the binary activation pattern $\boldsymbol{y}^{(S1)} \in \{0,1\}^{C^{(S1)}, H, W}$ as follows:

\begin{equation}
\boldsymbol{y}^{(S1)}_{c_{\text{out}}} = \Theta \left(\sum_{c_{\text{in}}=0}^C \boldsymbol{W}^{(S1)}_{c_{\text{out}},c_{\text{in}}} \star \boldsymbol{x}_{c_{\text{in}}} \right)
\label{eq:y_s1}
\end{equation}

Here, $\star$ denotes the convolution operator between the filters $\boldsymbol{W}^{(S1)} \in \mathbb{R}^{C^{(S1)}, C, h, w}$ and the input image $\boldsymbol{x}$, while the parameters  $h, w$ define the width and height of the convolutional filters, i.e., the spatial neighborhood from where the features are extracted.
To binarize the activations, we use the heaviside function $\Theta$, which sets a neuron to one if its input is larger than zero, and to zero otherwise.
%
\begin{equation}
  y_{i} = \Theta(\boldsymbol{a}_{i} ) = \begin{cases}
      1 \quad \text{if } \boldsymbol{a}_{i}  > 0\\
      0 \quad \text{otherwise}
  \end{cases}
  \label{eq:activation}
\end{equation}
%
Note that the width and height between input and output activations remain identical, but the number of channels changes from $C$ to $C^{(S1)}$, allowing the extraction of multiple features at each spatial location.

\subsection{Stage 2: Net Fragments}\label{sec:stage_2_impl}
$S2$ constructs net fragments from the binary feature activations $\boldsymbol{y}^{(S1)}$ of $S1$ and integrates them into coherent nets over time. To ensure stability, single neurons or small groups of neurons are suppressed through an attenuation mechanism, ensuring that only large-scale, consistent activation patterns persist. 

\paragraph{Network Topology.}
Stage 2 uses convolutional layers to implement the connections $w_{j,i}$ within net fragments $n_j \in \mathcal{F}(n_i)$.
Conceptually, this stage comprises a three-dimensional array of neurons with dimensions $C^{(S2)}\times H \times W$, where $C^{(S2)}$ denotes the number of feature channels in $S2$, and $H$ and $W$ correspond to the spatial dimensions of the input image.

The spatial distance between two neurons $n_i$ and $n_j$ is defined using the maximum norm (also known as Chebyshev distance), where adjacent neurons have a unit distance:

\begin{equation}
d(n_i,n_j) = \max \left( |h_i - h_j|, |w_i - w_j|\right)
\end{equation}

Each neuron $n^{(S2)}_i$ in stage 2 is connected in two ways: (1) \emph{Feedforward connections} to the corresponding local neighborhood $\mathcal{B}_r(n^{(S1)}_i)$ in $S1$, and (2) \emph{lateral connections} to neighboring neurons within stage 2 itself, $\mathcal{B}_r(n^{(S2)}_i)$, denoted as:

\begin{equation}
\begin{aligned}
    \mathcal{B}_r \left(n^{(S2)}_i \right) &= \mathcal{B}_r \left(n^{(S1)}_i \right) \cup \left\{n_j^{(S2)} \mid d(n_i^{(S2)}, n_j^{(S2)}) \leq r \right\}\\
    \mathcal{B}_r \left(n^{(S1)}_i \right) &= \left\{n_j^{(S1)} \mid d(n_i^{(S2)}, n_j^{(S1)}) \leq r \right\}
\end{aligned}
\end{equation}

where $r$ defines the radius of lateral connections. The convolutional architecture realizes these neighborhood structures via convolutional filters, where the filter size directly determines the radius of both feedforward and lateral connections.
Using convolutional filters not only implements the connectivity restrictions to local neighborhoods but also allows detection of the same pattern at various positions, enabling learning of net fragments independent of position.

The forward connections from $S1$ to $S2$, denoted as $\boldsymbol{W}^{(F)}$, are necessary to activate cells within $S2$. It can be interpreted as a small group of neurons in $S1$ (based on the filter size of $\boldsymbol{W}^{(F)}$) supporting a single neuron within $S2$.
In contrast, the lateral recurrent connections within $S2$, denoted as $\boldsymbol{W}^{(L)}$, are responsible for building the overlapping net fragments within $S2$, providing support amongst the neurons within $S2$.
It is important to note that the lateral connection matrix $\boldsymbol{W}^{(L)}$ does not implement recurrent processing in the traditional sense of sequential token-wise updates, as found in  RNN-based language models. Instead, $\boldsymbol{W}^{(L)}$ iterates over the same neuronal state vector $\boldsymbol{y}^{(S2)}$.
Specifically, the lateral connections $\boldsymbol{W}^{(L)}$ calculate the state updates of $\boldsymbol{y}^{(S2)}_{t\rightarrow t+1}$ within an iterative refinement process in which individual neuronal activations ${y}_i^{(S2)} \in \boldsymbol{y}^{(S2)}$ may flip due to changes in support and attenuation.

\paragraph{Pre-Synaptic Activity.}
To determine the support a single neuron $n_{i,t}$ receives from the net fragment $\mathcal{F}_{t-1}(n_{i, t-1})$, we first calculate the pre-synaptic activity $\mathcal{I}$ by applying convolutional operations on the forward and lateral connections and then taking the sum of these operations: 
%
\begin{equation}
\begin{aligned}
\boldsymbol{a}^{\mathcal{I}(S2)}_{c_{\text{out}}, t}
&= \underbrace{
    \sum_{c_\text{in}=0}^{C^{(S1)}} \boldsymbol{W}^{(F)}_{c_{\text{out}}, c_{\text{in}}} 
    \star \boldsymbol{y}_{c_{\text{in}}}^{(S1)}
}_{\textit{forward input } (\boldsymbol{a}^{\mathcal{I}(F)})}
+ 
\underbrace{
    \sum_{c_\text{in}=0}^{C^{(S2)}} \boldsymbol{W}^{(L)}_{c_{\text{out}}, c_{\text{in}}} 
    \star \boldsymbol{y}_{c_{\text{in}}, t-1}^{(S2)}
}_{\textit{lateral input } (\boldsymbol{a}^{\mathcal{I}(L)})}
\\[1.5ex]
\boldsymbol{a}^{\mathcal{I}(F)}_{c_{\text{out}}}
&= \sum_{c_\text{in}=0}^{C^{(S1)}} \sum_{h=0}^{H-1} \sum_{w=0}^{W-1} 
  \boldsymbol{W}^{(F)}_{c_{\text{out}}, c_{\text{in}}, h, w} \cdot 
  \boldsymbol{y}^{(S1)}_{c_{\text{in}}, h'+h, w'+w}
\\
\boldsymbol{a}^{\mathcal{I}(L)}_{c_{\text{out}}, t}
&= \sum_{c_\text{in}=0}^{C^{(S2)}} \sum_{h=0}^{H-1} \sum_{w=0}^{W-1} 
  \boldsymbol{W}^{(L)}_{c_{\text{out}}, c_{\text{in}}, h, w} \cdot 
  \boldsymbol{y}^{(S2)}_{c_{\text{in}}, h'+h, w'+w, t-1}
\end{aligned}
\label{eq:a_a2}
\end{equation}
%
Here, the forward connections have the shape $\boldsymbol{W}^{(F)} \in \mathbb{R}^{C^{(S2)}, C^{(S1)}, h, w}$ and the recurrent connections have the shape $\boldsymbol{W}^{(L)} \in \mathbb{R}^{C^{(S2)}, C^{(S2)}, h, w}$, where $C^{(S2)}$ is the number of channels in $S2$. 
Therefore, the pre-synaptic activity $\mathcal{I}$ at time $t$ is increased by the learned amount $\boldsymbol{w}_{j,i}$ from all connected and active neurons ($\boldsymbol{y}_{j,t-1} = 1$).
%
Higher support, indicated by larger $\boldsymbol{a}^{\mathcal{I}(S2)}_{i}$, suggests a more likely pattern, as neurons encoding frequent patterns gain reinforcement not only from feedforward connections but also from lateral connections.

Next, we extend this convolutional-based formulation of forward and lateral connections by introducing competitive neurons.
Afterward, we describe how the pre-synaptic activity $\boldsymbol{a}^{\mathcal{I}(S2)}$ undergoes attenuation, how it is transformed into $\boldsymbol{y}^{(S2)}$,
and how the weight matrices $\boldsymbol{W}^{(F)}, \boldsymbol{W}^{(L)}$ are initialized and trained.


\paragraph{Competitive Neurons.}
To address the challenge of context entanglement, wherein individual neurons may participate in multiple fragments, we implement \emph{competitive neurons} as follows: For each base feature channel $c \in C^{(S2)}$, we introduce $\kappa$ parallel alternatives that share the same receptive field structure but are allowed to develop distinct connections. Following our formal definition, each base neuron $n_{c,h',w'}$ at spatial location $(h',w')$ has the competitive class:

\begin{equation}
    \mathcal{C} \left( n_{c,h',w'} \right) = \left\{ n_{c,h',w'}^1, \dots, n_{c,h',w'}^\kappa\right\}
\end{equation}

We implement this by introducing an additional competition dimension of size $\kappa$ to the tensors in $S2$:
%
\begin{equation}
\begin{aligned}
   C^{(S2)} &\to (C^{(S2)},\,\kappa)\\
   \boldsymbol{W}^{(F)}\in\mathbb R^{C^{(S2)},\,C^{(S1)},\,h,\,w}
   &\to
   \mathbb R^{C^{(S2)},\,\kappa,\,C^{(S1)},\,h,\,w}\\
   \boldsymbol{W}^{(L)}\in\mathbb R^{C^{(S2)},\,C^{(S2)},\,h,\,w}
   &\to
   \mathbb R^{C^{(S2)},\,\kappa,\,C^{(S2)},\,h,\,w}\\
   \boldsymbol{a}^{\mathcal{I}(L)}\in\mathbb R^{C^{(S2)},H,W}
   &\to
   \mathbb R^{C^{(S2)},\,\kappa,\,H,W}
\end{aligned}
\end{equation}

Although this transformation increases the number of total feature maps by a factor of $\kappa$, only a single competitive neuron per base channel is active at any given spatial location (one neuron within $\mathcal{C} \left( n_{c,h',w'} \right)$). This sparsity is enforced through a group-wise \emph{winner-take-all (WTA)} mechanism along the $\kappa$-dimension.

At each time step $t$, competition operates independently within each class $\mathcal{C} \left( n_{c,h',w'} \right)$. For each base neuron, we compute pre-activation scores for all competitive neurons and select the winning competitive index:
%
\begin{equation}
  k^*(c,h',w') \;=\;
  \argmax_{k \in \{0,\dots,\kappa-1\}}\;
    \left\langle
      \boldsymbol{W}_{c,k,:,:,:,:},\;
      \boldsymbol{z}_{h',w',t}
    \right\rangle.
\end{equation}

where $\boldsymbol{z}_{h',w',t}$ is the concatenated input patch:

\begin{equation}
\boldsymbol{z}_{h',w',t} \;=\;
\begin{bmatrix}
\boldsymbol{y}^{(S1)}_{:,\, h':h'+H,\, w':w'+W} \\
\boldsymbol{y}^{(S2)}_{:,\, h':h'+H,\, w':w'+W,\, t-1}
\end{bmatrix}.
\end{equation}
%
Thus, at each spatial location and time step, the index $k^* \in \{1, \dots, \kappa \}$ maximizing the correlation between the corresponding filter and the local input patch is chosen.
Once the winning index $k^*$ is selected, the activation is computed exclusively for the corresponding neuron:
%
\begin{equation}
\begin{aligned}
\boldsymbol{a}^{\mathcal{I}(S2)}_{c, h', w', t} &=
\boldsymbol{a}^{\mathcal{I}(F)}_{c, k^*, h', w'} + \boldsymbol{a}^{\mathcal{I}(L)}_{c, k^*, h', w', t}
\\[1.5ex]
\boldsymbol{a}^{\mathcal{I}(F)}_{c, k^*, h', w'} &= 
\sum_{c_{\text{in}}=0}^{C^{(S1)} - 1} \sum_{h=0}^{H-1} \sum_{w=0}^{W-1}
\boldsymbol{W}^{(F)}_{c,k^*,\, c_{\text{in}},\, h,\, w} \cdot
\boldsymbol{y}^{(S1)}_{c_{\text{in}},\, h'+h,\, w'+w}
\\[1.5ex]
\boldsymbol{a}^{\mathcal{I}(L)}_{c, k^*, h', w', t} &= 
\sum_{c_{\text{in}}=0}^{C^{(S2)} - 1} \sum_{h=0}^{H-1} \sum_{w=0}^{W-1}
\boldsymbol{W}^{(L)}_{c,k^*,\, c_{\text{in}},\, h,\, w} \cdot
\boldsymbol{y}^{(S2)}_{c_{\text{in}},\, h'+h,\, w'+w,\, t-1}
\end{aligned}
\end{equation}

Initially, all competitive neurons $n_{c,h',w'}^k$ are identical due to shared initialization. However, Hebbian learning induces specialization: each competitive neuron within $\mathcal{C} \left( n_{c,h',w'} \right)$ gradually diverges, learning to encode different contextual variations of the same base feature.

In a practical implementation, we merge the $\kappa$ and channel dimensions, resulting in a flat channel axis of size $C^{(S2)} \cdot \kappa$. This allows the use of efficient standard 3D convolutional layers, where non-winning neurons are simply masked out at each spatial location by setting them to $0$, preserving the theoretical definition.


\paragraph{Attenuation.}
The pre-synaptic activity is processed by an activation function $\mathcal{A} \left(\boldsymbol{a}^\mathcal{I}, \mathcal{E} \right)$. In our implementation, this function first normalizes the potential  $\boldsymbol{a}^\mathcal{I} \rightarrow \boldsymbol{a}^{\text{norm}}$ to conform it into the range $[0, 1]$ (details see Appendix \ref{sec:normalization}), and then applies a globally increasing attenuation pressure $\boldsymbol{a}^{\text{norm}} \rightarrow \boldsymbol{a}$ to ensure that only neurons receiving substantial support remain active.
Specifically, attenuation is implemented as:
%
\begin{equation} 
    \boldsymbol{a}^{(S2)}_{c,t} = \left(\boldsymbol{a}^{\text{norm }(S2)}_{c,t}\right)^\gamma,
    \label{eq:inhibition}
\end{equation}
%
where $\gamma$ regulates the attenuation strength and increases over time according to $\gamma=\alpha+\beta t$ (when not stated differently, we set $\alpha=1.2$ and $\beta=0.2$), making activation progressively harder. This process leads to a sparsification of activations that occurs through both direct attenuation and reduced support from previously silenced neurons~.

The final output is then binarized using the heaviside function $\Theta$:
%
\begin{equation} 
     \boldsymbol{y}_{c,t}^{(S2)} = \Theta\left(\boldsymbol{a}^{(S2)}_{c,t}  - b^{(S2)} \right)
     \label{eq:s2_final_output}
\end{equation}
%
The activation at the last timestep $\boldsymbol{y}^{(S2)}_{t=T}$ represents a recurrently refined activation pattern where unsupported neurons have been suppressed, yielding a coherent representation.

\paragraph{Weight Initialization.}
The forward connections $\boldsymbol{W}^{(F)}$ are initialized so that the activations of the $S1$ neurons are initially copied into all competitive $S2$ neurons in the same position. This is done by setting the connections at the center of the kernel (at position $(h/2, w/2)$) to $1$ if $c_{in} = \lceil c_{out} / \kappa \rceil$, where $c_{in}$ denotes the index of the input channel, $c_{out}$ the index of the output channel, and $\kappa$ the number of competitive neurons. The recurrent connections $\boldsymbol{W}^{(L)}$ are initialized with zeros except for the self-coupling of neurons (recurrent connection to itself), which is set to $1$.

\paragraph{ Learning of Weights.}
Support within net fragments is \emph{learned} via Hebbian plasticity, so that frequently occurring patterns that lead to correlated activation patterns between neurons $n_i \sim n_j$ are ``stored'' in the system and receive more support:
%
\begin{equation}
    \boldsymbol{W}^{(S2)} := \min \left( \max \left( \boldsymbol{W}^{(S2)} + \eta \cdot \boldsymbol{\rho}_{avg}, 0\right), 1\right)
    \label{eq:hebbian_plasticity}
\end{equation}
%
Here, $\eta$ represents the learning rate and $\boldsymbol{W}^{(S2)}$ stands for both $\boldsymbol{W}^{(F)}$ and $\boldsymbol{W}^{(L)}$.
Note that the same convolutional kernel is applied at various spatial positions. The decision to update this connection is based on the average correlation $\boldsymbol{\rho}_{avg}$ between all input and output neurons it connects (across all spatial positions).
If the resulting average is positive, the corresponding connection increases as it connects more simultaneously active neurons than neurons that fire disjointly. Conversely, if the average correlation is negative, indicating more disjoint firing, the connection strength is reduced.
Importantly, the weight update is applied only once per sample, after the full sequence of $T$ timesteps has been processed. Learning proceeds over multiple epochs, meaning each sample is revisited multiple times. Due to the small learning rate of $\eta=0.2$, these weight changes accumulate gradually.
